% Generated draft: requires mathematical and source review.
\documentclass[11pt]{article}
\usepackage[margin=25mm]{geometry}
\usepackage{fontspec}
\setmainfont{Latin Modern Roman}
\usepackage{amsmath,amssymb,mathtools,unicode-math}
\setmathfont{Latin Modern Math}
\usepackage{enumitem,needspace,xurl,hyperref}
\setlength{\parindent}{0pt}
\setlength{\parskip}{5pt}
\setlength{\emergencystretch}{3em}
\newcommand{\N}{\mathbb{N}}
\newcommand{\Z}{\mathbb{Z}}
\newcommand{\Q}{\mathbb{Q}}
\newcommand{\R}{\mathbb{R}}
\newcommand{\C}{\mathbb{C}}
\newcommand{\F}{\mathbb{F}}
\newcommand{\A}{\mathbb{A}}
\newcommand{\PP}{\mathbb{P}}
\newcommand{\E}{\mathbb{E}}
\newcommand{\eps}{\varepsilon}
\newcommand{\dd}{\,\mathrm{d}}
\begin{document}
\section*{UnsolvedMath ID: 30005361}
\textbf{Source problem code:} catalogue entry
\subsection{Definitions and mathematical statement}
Let \(k\geq 1\) be an integer. Let \(\mathcal K\) be a first-order structure with a distinguished field \(K\), and put \(T=\operatorname{Th}(\mathcal K)\), in the full language of that structure. For a pure field, this is the language of rings; for a field with additional structure, all that structure is included when testing the following properties.

A formula \(\varphi(x;y_1,\ldots,y_k)\), whose displayed variables may be finite tuples, has the \(k\)-independence property in \(T\) if there exist a model \(M\models T\), parameter tuples \(a_i^j\) for \(1\leq j\leq k\) and \(i\in\N\), and tuples \(b_S\) for every subset \(S\subseteq\N^k\), such that
\[
M\models\varphi(b_S;a_{i_1}^1,\ldots,a_{i_k}^k)
\quad\Longleftrightarrow\quad
(i_1,\ldots,i_k)\in S
\]
for every \(S\subseteq\N^k\) and every \((i_1,\ldots,i_k)\in\N^k\). The theory \(T\) is \(k\)-dependent if no formula has this property. It is \(\mathrm{NIP}\) if it is \(1\)-dependent. The structure \(\mathcal K\) is called \(k\)-dependent or \(\mathrm{NIP}\) according to its complete theory.

The supplied Conjecture 1 asks whether, for every positive integer \(k\) and every field structure \(\mathcal K\) within its intended scope,
\[
\operatorname{Th}(\mathcal K)\text{ is }k\text{-dependent}
\quad\Longrightarrow\quad
\operatorname{Th}(\mathcal K)\text{ is }\mathrm{NIP}.
\]
In particular, the pure-field assertion has the precise scope
\[
\forall k\in\Z_{\geq1}\;\forall K\text{ a field},\qquad
\operatorname{Th}(K;+,-,\cdot,0,1)\text{ is }k\text{-dependent}
\ \Longrightarrow\ 
\operatorname{Th}(K;+,-,\cdot,0,1)\text{ is }\mathrm{NIP}.
\]
No restriction on characteristic, cardinality, or finiteness of the field is supplied. For \(k=1\), the implication holds by the definitions; the substantive cases are \(k\geq2\).

The source explicitly includes fields equipped with a valuation or a derivation, and also says ``etc.'' The implication for such an expansion concerns its full theory, including the additional structure. The source does not specify the complete class of permitted expansions, their languages or sorts, or additional compatibility hypotheses. Consequently, the expanded-field assertion above is a schema for the intended field structures, rather than a fully specified universal assertion over all possible expansions.

As motivation, the extracted passage reports that \(k\)-dependent fields are Artin--Schreier closed, citing its reference [7] and describing this as a generalization of Kaplan--Scanlon--Wagner for \(k=1\). It also reports Henselianity for valued fields of positive characteristic in the surrounding \(k\)-dependent setting, citing [2] and describing this as a generalization of Johnson for \(k=1\), and attributes a reduction of the valued-field question to the pure-field question to Boissonneau. These are supplied background claims, not additional hypotheses or independently verified results in this draft; their complete hypotheses and the identities of [7] and [2] are unavailable in the record.

The extraction additionally contains a separate speculation that a \(k\)-dependent theory should be ``linear, or \(1\)-based'' relative to its \(\mathrm{NIP}\) part, together with incomplete references to Theorem 1 and a missing precise version. That fragment does not supply enough definitions or text to formulate an additional conjecture here.
\subsection{Short English statement}
The conjecture asks whether every \(k\)-dependent field is already \(\mathrm{NIP}\), for every positive integer \(k\). The source includes pure fields and fields with valuations, derivations, or other unspecified additional structure. Dependence is tested in the full language of the structure. The pure-field statement is precise; the full range of intended expansions remains unresolved. The supplied record labels the problem open, subject to its dated, unverified literature assessment.
\subsection{Sources}
\begin{itemize}
\item Local UnsolvedMath problem registry, permanent ID \(30005361\), title: NIP Conjecture for \(k\)-Dependent Fields; exact source problem code: \texttt{catalogue entry}. This draft uses the supplied complete source record. The record assigns the problem to the Oberwolfach Reports Open Problems collection, gives proposed year 2023, and supplies no individual proposer.
\item Supplied original report citation: \emph{Model Theory: Combinatorics, Groups, Valued Fields and Neostability} (2023), Oberwolfach Reports, pp. 91--149, DOI 10.4171/owr/2023/2. The extracted statement is labeled Conjecture 1. \url{https://doi.org/10.4171/owr/2023/2}. This is also the supplied literature-source link labeled ``original Oberwolfach report.''
\item Supplied literature source: Chernikov--Hempel, \emph{n-dependent fields II}, linked to \emph{On n-dependent groups and fields II}. \url{https://www.cambridge.org/core/journals/forum-of-mathematics-sigma/article/on-ndependent-groups-and-fields-ii/0C88DB70DC29AAD54EFE858F1679D011}.
\item Supplied literature source: Chernikov--Hempel, \emph{n-dependent fields III}. \url{https://arxiv.org/abs/2412.19921}.
\item Literature assessment dated exactly \texttt{2026-08-22}: the supplied curation labels the conjecture open and reports structural consequences and partial results, but no proof or counterexample to the general assertion. This is supplied dated curation, not an independently verified literature assessment. The background expressly calls for independent verification of time-sensitive status and cited claims.
\item The extracted background attributes earlier \(k=1\) results to Kaplan--Scanlon--Wagner and Johnson, and a reduction concerning valued fields to Boissonneau. Its reference labels [7] and [2] are retained as supplied; their bibliographic entries are absent from this record.
\end{itemize}
\end{document}
